\chapter{Conclusão}
\label{cha:conclusao}

Este trabalho analisou a literatura científica sobre o metaverso a partir de diferentes perspectivas: definição, tecnologias, arquiteturas, aplicações e interoperabilidade. A análise foi deliberadamente limitada a informações que puderam ser verificadas nas próprias publicações utilizadas como fonte.

Os resultados mostram que não existe uma definição única consolidada. Ritterbusch e Teichmann analisaram 28 definições e descrições; Boztas et al. analisaram 103 estudos, dos quais 33 apresentavam definições explícitas e 25 apresentavam modelos arquiteturais; Tukur et al. analisaram técnicas, tecnologias e aplicações a partir de 30 estudos primários; e Enamorado-Díaz et al. selecionaram 35 estudos primários sobre aplicações e desafios de engenharia de software. \cite{ritterbusch2023,boztas2025,tukur2024,enamorado2025}

A comparação permite sustentar uma definição operacional em camadas. O primeiro nível descreve o fenômeno e suas propriedades; o segundo descreve as tecnologias que podem sustentá-lo; e o terceiro descreve suas aplicações. Essa estrutura reduz o risco de definir o metaverso por uma tecnologia específica.

A principal correção metodológica em relação às versões anteriores foi retirar números que não correspondiam a uma busca própria documentada. O presente trabalho não apresenta um PRISMA numérico próprio porque não dispõe de uma cadeia completa e auditável de identificação, duplicação, triagem e elegibilidade executada pelo estudo. Em seu lugar, os quantitativos das revisões analisadas são apresentados com sua fonte explícita.

Essa escolha é importante para a integridade acadêmica do trabalho. PRISMA 2020 orienta o relato transparente de revisões sistemáticas, enquanto PRISMA-S detalha a documentação das buscas para que possam ser reproduzidas. \cite{page2021,rethlefsen2021}

A definição operacional proposta é, portanto, uma síntese para este projeto, e não uma afirmação de consenso científico. Uma etapa futura poderá transformar o estudo em uma revisão sistemática própria, desde que sejam executadas as strings nas bases selecionadas, registrados os resultados de cada base, exportados os registros, removidas as duplicatas, documentados os motivos de exclusão e construído o fluxograma PRISMA a partir desses dados.

Até que essa etapa seja efetivamente realizada, qualquer número adicional de registros da revisão própria deve ser considerado inexistente, e não estimado.
